\section{The inert process has the original mean exponent}
\label{gs:inert-exponent-section}

Let $\mathcal X$ be the class of finite birth environments, allowing
inert boundary stubs, with the tag initially satisfying $L_u=c_u=1$.
For $x\in\mathcal X$, let $Z_m(t;x)$ count its future core descendants
of genealogical generation $m$ born by time $t$, and put
\[
 F_x(t)=1+\sum_{m\ge1}\mathbb E_x Z_m(t;x).
\]
The initial tag contributes the term $1$; other initial core vertices
are not counted. Write $x_\infty$ for the tag with one leaf and one
inert parent stub, and $F_\infty=F_{x_\infty}$.

The input from~\cite{Uzay} is the original free-tree mean exponent
\[
 F_2(t):=\mathbb E_{\mathrm{edge}} C(t)
       =e^{\mu t+o(t)}.
\]
We also use the planted-branch comparison proved in its
\texttt{proofs/existence.md} in Lemma 4.1(c), Corollary 5.3, and Lemma 6.1 of that source. We state precisely the portion used here.
Fix $M\ge2$, sufficiently small $\epsilon>0$, and
$k\ge6M/\epsilon$. Put
\[
 \eta=b(\epsilon)\log2+\frac9M,
 \qquad b(\epsilon)=7/\log(1/\epsilon)\longrightarrow0
       \quad(\epsilon\downarrow0).
\]
For all sufficiently large $t$, there is a rooted $k$-vertex free tree
$(T,v)$ for which the following lower comparison holds. If $T$ is
planted as a branch with its root $v$ attached to an exterior vertex
of full degree at least $M$, its future expected internal core count
is at least
\begin{equation}\label{gs:planting-input}
 \frac1k e^{-\eta t}
       \bigl[F_2(t)-k-O(e^{-t})\bigr].
\end{equation}
The exterior may evolve freely. The constants are uniform over the
finite collection of rooted $k$-vertex seeds. The seed and root can
depend on $t$, but are selected from their laws before planting.

For clarity, the comparison weights a free reference history by
\[
 e^{-t/M}\prod_{\tau\le t:\,\text{source }v}
       \frac{d_v(\tau-)}{d_v(\tau-)+1},
\]
where $d_v$ is its internal reference degree. A high-degree original
vertex in the reference history gives the factor $e^{-\eta t}$;
averaging over its $k$ possible choices gives $1/k$. The inequality
$F_2(t)\le k+\max_{|T|=k}\mathbb E C_T(t)$ supplies the displayed
free-tree term. This is the cited comparison input, not a new
independence assumption about the two sides of an edge.

\begin{lemma}[Uniform finite preparation]\label{gs:uniform-preparation}
For fixed $M,k$, there are $h<\infty$ and $p>0$ such that, from every
$x\in\mathcal X$, any chosen rooted $k$-vertex seed $(T,v)$ can be
planted by time $h$ with probability at least $p$. Every core vertex in
the branch belongs to the original tag's future core genealogy, and its
root is attached to that tag, whose full degree is then at least $M$.
Both constants are uniform
in the exterior environment and in the rooted seed.
\end{lemma}

\begin{proof}
If the tag has a real parent, reserve potential parent-to-tag arrows
at rate $1/2$ and accept each with probability $2/D_p(t)$. The parent's
full degree is at least two and never decreases, so this is an exact
construction. Suppressing these potential arrows during a fixed
preparation interval costs only $e^{-h/2}$, independently of the
retained component's clocks. For a stub parent the same suppression
can be imposed on an independent unused stream. The preparation can
therefore be performed in the inert-parent process.

First keep the tag silent for one unit of time and let its one leaf
produce enough Yule offspring that the tag's full degree is at least
$M$. This has positive probability. In a second unit interval require
exactly one tag ring, choosing a leaf, and keep the resulting two-vertex
branch silent until that interval ends. The leaf-choice probability is
at least $(M-1)/M$, since all growth before that promotion increases
the number of leaves at the tag. Its full degree is unchanged by the
promotion.

Starting with that branch edge, build the specified rooted tree by
adjoining its remaining vertices in an order in which each new vertex
has an already present parent. During each unit interval keep the tag
silent and prescribe one ring of an internal neighbour of the desired
attachment vertex, its choice, and silence of all other internal clocks.
Also require the newborn to remain silent until the interval ends.
All relevant internal degrees are at most $k+1$, so every prescribed
choice has probability at least $1/(k+1)$. Old leaves outside the branch
can grow freely; they only increase the exterior tag's degree.

There are finitely many prescribed clocks and choices and finitely many
rooted seeds of size $k$. Taking their minimum probability gives
$p>0$, with a common preparation time, for example $h=k$.
\end{proof}

\begin{theorem}[Uniform lower growth and the inert exponent]
\label{gs:inert-mu}
For every $\varepsilon>0$ there are $c_\varepsilon>0$ and
$t_\varepsilon<\infty$ such that
\begin{equation}\label{gs:uniform-lower-time}
 \inf_{x\in\mathcal X}F_x(t)
       \ge c_\varepsilon e^{(\mu-\varepsilon)t}
       \qquad(t\ge t_\varepsilon).
\end{equation}
Moreover the inert process has the exact mean exponent
\begin{equation}\label{gs:inert-mean-limit}
 \lim_{t\to\infty}\frac{\log F_\infty(t)}t=\mu.
\end{equation}
\end{theorem}

\begin{proof}
After the uniform preparation, apply~\eqref{gs:planting-input} to its
fresh branch. Every reference core counted by that comparison is an
actual future descendant of the original tag. Thus, for all large $t$,
\begin{equation}\label{gs:uniform-planted-bound}
 \inf_{x\in\mathcal X}F_x(h+t)
 \ge\frac pk e^{-\eta t}
          \bigl[F_2(t)-k-O(e^{-t})\bigr].
\end{equation}
Uniformity of $p,h$ over the finite seed set permits the optimizing
seed to vary with $t$. Fix the comparison parameters before taking
$t\to\infty$, and then let $\eta\downarrow0$. This proves
\eqref{gs:uniform-lower-time}, and in particular the required lower
bound for the exponent of $F_\infty$.

For the upper bound, embed the inert tag and its leaf in a finite real
tree $T_M$ whose parent has full degree $M$: give the parent $M-1$
additional leaves. Reserve its potential arrow to the tag at rate
$1/M$, accepting with probability $M/D_p(t)$. Before the first
potential arrow, the retained tag component follows exactly the inert
law; a tag ring choosing the exterior affects only the parent and
does not change the retained component. The potential-arrow stream is
independent of the inert component's drivers. Therefore
\begin{equation}\label{gs:inert-upper-embedding}
 e^{-t/M}F_\infty(t)\le\mathbb E C_{T_M}(t).
\end{equation}
Every fixed finite tree is reachable from the original edge by a fixed
time with positive probability. The Markov property consequently gives
constants $a_M,b_M<\infty$ with
\[
 \mathbb E C_{T_M}(t)\le a_MF_2(t+b_M).
\]
For fixed $M$,~\eqref{gs:inert-upper-embedding} implies
$\limsup t^{-1}\log F_\infty(t)\le\mu+1/M$. Letting $M\to\infty$
finishes the proof.
\end{proof}

